% =====================================================================
%  Documentos largos (informe / tesina) — suite de documentos LaTeX
%  Nivel 3 (Tipos de documento) · clase report · autodemostrativo.
%  Compilar: latexmk -pdf documentos-largos.tex
% =====================================================================
\documentclass[11pt,a4paper]{report}

% --- Base estándar ---
\usepackage[utf8]{inputenc}
\usepackage[T1]{fontenc}
\usepackage{lmodern}
\usepackage{microtype}
\usepackage[spanish,es-noshorthands]{babel}
\usepackage{csquotes}
\usepackage{amsmath}\usepackage{amssymb}
\usepackage{booktabs}
\usepackage[margin=2.5cm]{geometry}
\usepackage{enumitem}
\usepackage{xcolor}
\usepackage{listings}
% --- Específicos de documentos largos ---
\usepackage{fancyhdr}      % encabezados/pies
\usepackage{setspace}      % interlineado
\usepackage{makeidx}\makeindex   % índice analítico
\usepackage{hyperref}
\hypersetup{unicode, pdftitle={Documentos largos en LaTeX},
            pdfauthor={Rodrigo Martín Vidal Galán},
            colorlinks=true, linkcolor=blue!50!black, urlcolor=blue!50!black}
\usepackage[spanish,capitalise,nameinlink]{cleveref}

% --- Estilo de código ---
\definecolor{codeback}{RGB}{245,247,250}
\definecolor{codekw}{RGB}{20,60,120}
\definecolor{codecom}{RGB}{110,120,135}
\lstset{
  basicstyle=\ttfamily\footnotesize,
  keywordstyle=\color{codekw}\bfseries,
  commentstyle=\color{codecom}\itshape,
  backgroundcolor=\color{codeback},
  frame=single, rulecolor=\color{codekw!30},
  breaklines=true, showstringspaces=false, columns=fullflexible,
  language=[LaTeX]TeX, xleftmargin=8pt, framesep=4pt, extendedchars=true,
  literate={á}{{\'a}}1 {é}{{\'e}}1 {í}{{\'i}}1 {ó}{{\'o}}1 {ú}{{\'u}}1
           {ñ}{{\~n}}1 {¿}{{?`}}1 {¡}{{!`}}1,
}
\newcommand{\C}[1]{\texttt{\textbackslash #1}}
\newcommand{\pkg}[1]{\textsf{#1}}

% --- Encabezados/pies con fancyhdr (demostración) ---
\pagestyle{fancy}
\fancyhf{}
\fancyhead[L]{\nouppercase{\leftmark}}  % capítulo a la izquierda
\fancyhead[R]{\thepage}                  % página a la derecha
\renewcommand{\headrulewidth}{0.4pt}

\title{\textbf{Documentos largos en \LaTeX}\\[4pt]
       \large Informe / tesina: estructura, numeración y modularización}
\author{Rodrigo Martín Vidal Galán}
\date{Junio de 2026}

\begin{document}

% --- Carátula ---
\maketitle
\thispagestyle{empty}

% --- Índice(s) en hojas aparte ---
\tableofcontents
\listoffigures
\listoftables
\clearpage

% =====================================================================
\chapter{Estructura jerárquica}
% =====================================================================
Los documentos largos usan las clases \texttt{report} o \texttt{book} (o sus
versiones KOMA, \texttt{scrreprt}/\texttt{scrbook}), que habilitan el nivel
\C{chapter}. La jerarquía completa, de mayor a menor:
\begin{center}\small
\renewcommand{\arraystretch}{1.2}
\begin{tabular}{@{}lll@{}}
\toprule
\textbf{Comando} & \textbf{Nivel} & \textbf{Numeración} \\
\midrule
\C{part}          & $-1$ & Parte I, II\dots \\
\C{chapter}       & $0$  & 1, 2, 3 (solo report/book) \\
\C{section}       & $1$  & 1.1, 1.2 \\
\C{subsection}    & $2$  & 1.1.1 \\
\C{subsubsection} & $3$  & 1.1.1.1 \\
\C{paragraph}     & $4$  & (sin número) \\
\bottomrule
\end{tabular}
\end{center}

\noindent La versión con \texttt{*} (p.\,ej.\ \C{section*\{\dots\}}) crea el
título \textbf{sin numerar} y sin entrada en el índice; para que igual aparezca
en el índice se usa \C{addcontentsline}.

\section{Front, main y back matter}
En \texttt{book} (y KOMA) se separan tres zonas con \C{frontmatter} (portada,
índices; numeración romana \texttt{i, ii}), \C{mainmatter} (el cuerpo;
reinicia a \texttt{1}) y \C{backmatter} (apéndices, bibliografía, índice). En
\texttt{report} se imita con \C{pagenumbering}:
\begin{lstlisting}
\pagenumbering{roman}   % i, ii, iii para indices
\tableofcontents
\pagenumbering{arabic}  % 1, 2, 3 desde el cuerpo
\end{lstlisting}

\section{Una sección de ejemplo}\label{sec:ejemplo}
Este texto sirve para que el \cref{cap:numeracion} muestre referencias cruzadas
\textbf{entre capítulos}: \C{cref} resuelve el número aunque esté en otro
archivo. Véase también la \cref{fig:demo} y la \cref{tab:demo}.

\begin{figure}[h]
  \centering
  \fbox{\rule{0pt}{2cm}\rule{4cm}{0pt}}
  \caption{Figura de ejemplo (entra en el \emph{Índice de figuras}).}
  \label{fig:demo}
\end{figure}

\begin{table}[h]
  \centering
  \begin{tabular}{@{}lr@{}}\toprule Concepto & Valor \\\midrule $a$ & 1 \\ $b$ & 2 \\\bottomrule\end{tabular}
  \caption{Tabla de ejemplo (entra en el \emph{Índice de tablas}).}
  \label{tab:demo}
\end{table}

% =====================================================================
\chapter{Numeración y profundidad}\label{cap:numeracion}
% =====================================================================
\LaTeX{} numera solo, pero podés controlar \textbf{cuánto} numera y muestra:
\begin{lstlisting}
\setcounter{secnumdepth}{2}  % numera hasta subsection (nivel 2)
\setcounter{tocdepth}{2}     % el indice muestra hasta subsection
\end{lstlisting}
Subir \texttt{secnumdepth} numera más niveles; bajarlo, menos. \texttt{tocdepth}
hace lo mismo para el índice.

\section{Atar la numeración a los capítulos}
Por defecto, en \texttt{report}, las ecuaciones se numeran corrido (1, 2, 3\dots)
en todo el documento. Para que se numeren \textbf{por capítulo} (1.1, 1.2,
2.1\dots) se usa \C{numberwithin} de \pkg{amsmath}:
\begin{lstlisting}
\numberwithin{equation}{chapter}   % ecuaciones: 1.1, 1.2, 2.1...
\numberwithin{figure}{chapter}
\end{lstlisting}
El mismo mecanismo reinicia el contador al empezar cada capítulo. Para forzar un
valor puntual: \C{setcounter\{equation\}\{0\}}.

\section{Reiniciar y dar formato a páginas}
\C{pagenumbering\{roman\}}/\texttt{\{arabic\}} cambia el estilo \emph{y} reinicia
el contador. Para el estilo de página (qué va en encabezado/pie) se usa
\C{pagestyle}; este documento define uno con \pkg{fancyhdr}:
\begin{lstlisting}
\pagestyle{fancy}\fancyhf{}
\fancyhead[L]{\nouppercase{\leftmark}}  % nombre del capitulo
\fancyhead[R]{\thepage}
\renewcommand{\headrulewidth}{0.4pt}
\end{lstlisting}
\C{leftmark} contiene el capítulo actual (mirá el encabezado de esta página).

% =====================================================================
\chapter{Modularización}
% =====================================================================
Un documento de 100 páginas en un solo \texttt{.tex} es inmanejable. La práctica
es \textbf{partirlo en archivos} y un \enquote{maestro} que los reúne:
\begin{lstlisting}
% archivo maestro: tesina.tex
\documentclass{report}
... preambulo ...
\begin{document}
  \include{cap1-intro}     % cada \include EMPIEZA en pagina nueva
  \include{cap2-metodos}
  \include{cap3-resultados}
\end{document}
\end{lstlisting}

\section{\texttt{\textbackslash input} vs.\ \texttt{\textbackslash include}}
\begin{itemize}[leftmargin=*]
  \item \C{input\{archivo\}} \enquote{pega} el contenido tal cual, en cualquier
        lado (sirve para el preámbulo o trozos pequeños).
  \item \C{include\{archivo\}} agrega un \C{clearpage} antes y después (pensado
        para \textbf{capítulos}) y, sobre todo, habilita\dots
  \item \C{includeonly\{cap2-metodos\}} en el preámbulo: compila \textbf{solo}
        ese capítulo, pero \textbf{conservando} la numeración y las referencias
        del resto. Acelera muchísimo mientras trabajás en una parte.
\end{itemize}
La alternativa moderna es el paquete \pkg{subfiles}, donde cada capítulo es
compilable \textbf{por sí solo} (con su propio preámbulo compartido).

\section{Buenas prácticas de organización}
\begin{itemize}[noitemsep,leftmargin=*]
  \item Preámbulo en un archivo aparte (\C{input\{preambulo\}}); reutilizable.
  \item Un archivo por capítulo, nombrado con sentido.
  \item Figuras en \texttt{figuras/}; bibliografía en \texttt{refs.bib}.
  \item \textbf{Una oración por línea} en el fuente: diffs de git limpios.
\end{itemize}

% =====================================================================
\appendix
\chapter{Apéndice: índices especiales}
% =====================================================================
\C{appendix} cambia la numeración de capítulos a letras (A, B\dots). Acá van
materiales complementarios.

Un \textbf{índice analítico} (\emph{index}) se arma con \pkg{makeidx}: marcás
términos con \C{index\{término\}}\index{índice analítico}\index{makeidx} y al
final \C{printindex} los lista con sus páginas. Requiere una corrida de
\texttt{makeindex} (que \texttt{latexmk} hace solo).

\printindex
\addcontentsline{toc}{chapter}{Índice analítico}

\end{document}
